\documentclass[10pt]{article}
\usepackage[margin=1in]{geometry}

% phfnote is a package for short documentation
% silence the warning associated with phfnote
\usepackage{silence}
\WarningFilter{latex}{Command \underbar  has changed.}
\WarningFilter{latex}{Command \underline  has changed.}
\usepackage[title=small]{phfnote}

\usepackage{amsmath}
\usepackage{amssymb}
\usepackage{amsthm}
\usepackage{bm}
\usepackage{bbm}
\usepackage{mathtools}
\allowdisplaybreaks
\usepackage{multirow}

\usepackage{graphicx}
\usepackage{caption}
\usepackage{subcaption}
\usepackage{float}
\usepackage{hyperref}
\usepackage{enumitem}
\usepackage{booktabs}
\usepackage{natbib}

\DeclareMathOperator*{\plim}{plim}
\DeclareMathOperator*{\var}{Var}
\DeclareMathOperator*{\cov}{Cov}
\DeclareMathOperator*{\corr}{corr}
\DeclareMathOperator*{\rank}{rank}
\DeclareMathOperator*{\tr}{tr}

\newcommand{\indep}{\perp \!\!\! \perp}
\newcommand{\R}{\mathbb{R}}
\newcommand{\I}{\mathbbm1}
\newcommand{\argmax}{\mathop{\rm arg~max}\limits}
\newcommand{\argmin}{\mathop{\rm arg~min}\limits}

% Theorem environments
\theoremstyle{plain}
\newtheorem{theorem}{Theorem}
\newtheorem{corollary}{Corollary}
\newtheorem{lemma}{Lemma}
\newtheorem{proposition}{Proposition}
\theoremstyle{definition}
\newtheorem{definition}{Definition}
\newtheorem{assumption}{Assumption}
\newtheorem{remark}{Remark}

\title{Reparameterization of the utility function}
\author{Konan Hara}
\date{\today}

\begin{document}

\maketitle

This note proposes a replacement for Section~\ref{sec:model} of the paper. It
takes total project net present value (NPV) as the primitive payoff and
distinguishes the bidder's absolute risk aversion over total NPV, $A_i$, from
the effective risk-aversion coefficient governing choices expressed per unit
of lifetime production, $\gamma_i=A_iK_i$. Here,
$K_i\equiv Capacity_iHT$ denotes lifetime production. This reparameterization
leaves the equilibrium bid-share equations, the sealed-bid cost inversion,
and the descending-clock solution unchanged when written in terms of
$\gamma_i$.

\section{Structural Model of New Energy Auctions}
\label{sec:model}

I model bidders participating in a multi-unit procurement auction following
the share auction framework of \citet{Wilson1979}. The distinguishing feature
of the model is that bidders bid a share of production that they will include
in the awarded long-term purchase agreement. Risk-averse bidders essentially
make portfolio choices to allocate their total production across the
risk-free purchase agreement and the risky wholesale market.\footnote{\label{fn:uncertainty}
The production amount and purchase agreement price can be uncertain at the
time of bidding due to intermittency and counterparty risk, respectively. The
model can incorporate these uncertainties by specifying their distribution.
The bidder's risk premium will then represent the risk premium for the
wholesale market relative to the purchase agreement, and the bidder's cost
will reflect the price he expects to receive, as in \citet{Ryan2026}.
} Each bidder also bids one price per unit, which applies to all units of the
purchase agreement.\footnote{
The model is a special case of a multi-unit auction model because one price is
applied to all units rather than bidding a price schedule along each share.
}

An auctioneer holds procurement auctions that guarantee the purchase of
electricity at a fixed price for the entire life of a renewable technology,
$T$. An auction at time $t=0$ is characterized by a lead time $l$, a number of
participants $N$, a procurement capacity $D$, and a minimum possible bid
share $\underline{q}\in[0,1]$. Qualified bidders, $i=1,\dots,N$, each with a
new investment project, compete for the procurement capacity $D$. The
procurement capacity is not disclosed before bidding, which makes it a random
variable from the bidders' perspective. Bidders are required to allocate at
least a share $\underline{q}$ of their total production to the purchase
agreement.

The purchase agreement spans discrete time
$t=l,l+1,\dots,l+T-1$, since electricity supply begins at time $t=l$ and
lasts for $T$ periods. Each period consists of $H$ hours. Bidder $i$ has
generation capacity $Capacity_i$ and therefore produces
$Capacity_iH$ units of electricity per period. Define
\begin{align*}
K_i\equiv Capacity_iHT
\end{align*}
as the project's lifetime production. Bidder $i$ specifies a share
$q_i\in[\underline{q},1]$ and a price $b_i$ in his bid.\footnote{
In practice, because bid quantities must be submitted in discrete units
(integers in the first auction in 2009 and up to the first decimal place in
subsequent auctions), the exact upper and lower bounds can deviate slightly
from the nominal $\underline{q}$ and $1$ and vary at the project level. I
maintain continuous, uniform bounds in the theoretical model for clarity and
address these discrete project-level constraints in the estimation.
} If bidder $i$ wins, the auctioneer agrees to purchase
$q_iCapacity_iH$ units of electricity per period at price $p_i$, which is
determined as a function of participants' bid prices
$b_1,\dots,b_N$ according to the auction format. Bidder $i$ sells the
remaining production, $(1-q_i)Capacity_iH$, to the wholesale market at price
$r_t$. The auctioneer awards purchase agreements to the lowest bid prices
until the total bid capacity of winners exceeds $D$. Thus, bidder $i$ wins
when
\begin{align*}
\sum_{j\neq i}q_jCapacity_j\I(b_j\leq b_i)<D.
\end{align*}

When bidder $i$ wins, he incurs an upfront fixed cost $FC_i$ and variable cost
$VC_{it}$ per unit of production in period $t$. His total project NPV is
\begin{align}
\label{eq:totalnpv}
\Pi_i
&\equiv
\sum_{t=l}^{l+T-1}Capacity_iH\delta^t
\left\{q_ip_i+(1-q_i)r_t-VC_{it}\right\}-FC_i \notag\\
&=K_i\left\{q_i\widetilde{\delta}p_i+(1-q_i)r-c_i\right\},
\end{align}
where $\delta$ is a common discount factor,
\begin{align*}
\widetilde{\delta}
\equiv \frac{1}{T}\sum_{t=l}^{l+T-1}\delta^t,
\qquad
r\equiv \frac{1}{T}\sum_{t=l}^{l+T-1}\delta^tr_t,
\end{align*}
and
\begin{align*}
c_i\equiv
\frac{FC_i}{K_i}
+\frac{1}{T}\sum_{t=l}^{l+T-1}\delta^tVC_{it}
\end{align*}
is cost per unit of lifetime production. Although preferences are defined over
total NPV, retaining $c_i$ in per-unit terms makes bid prices, costs, risk
premiums, and markups directly comparable in dollars per unit of electricity.

I assume that bidder $i$ has constant absolute risk aversion (CARA) over total
project NPV:
\begin{align*}
U_i(\Pi\mid A_i)
=\frac{1-\exp(-A_i\Pi)}{A_i},
\qquad A_i>0,
\end{align*}
where $A_i$ is the Arrow--Pratt coefficient of absolute risk aversion over
total NPV. Exponential utility is defined over the entire real line and is
therefore compatible with projects having negative realizations of NPV. A
concave utility over project NPV is a viable option for analyzing projects
with uncertain cash flows in the field of decision analysis
\citep{Hazen2009}.

Define the effective risk-aversion coefficient
\begin{align}
\label{eq:effectivegamma}
\gamma_i\equiv A_iK_i=A_iCapacity_iHT.
\end{align}
If $\pi_i\equiv\Pi_i/K_i$ denotes NPV per unit of lifetime production, then
\begin{align}
\label{eq:utilityequivalence}
U_i(K_i\pi_i\mid A_i)
=K_i u(\pi_i\mid\gamma_i),
\qquad
u(\pi\mid\gamma_i)
\equiv\frac{1-\exp(-\gamma_i\pi)}{\gamma_i}.
\end{align}
For a bidder with predetermined capacity, $K_i$ is a positive constant that
does not vary with the bid $(q_i,b_i)$. Hence, maximizing expected utility over
total NPV is equivalent to maximizing the normalized per-unit utility on the
right-hand side of equation~\eqref{eq:utilityequivalence}. Capacity affects
the bidder's exposure mechanically through $K_i$, while $A_i$ measures the
bidder's primitive attitude toward total project risk. Portfolio choices
identify their product $\gamma_i=A_iK_i$; because $K_i$ is observed, they also
identify $A_i$ under the model.

I allow absolute risk aversion over total NPV to vary systematically with
project capacity as
\begin{align}
\label{eq:Acapacity}
A_i=\exp\!\left(\beta_{A0}+\beta_{A1}\log Capacity_i\right).
\end{align}
Equations~\eqref{eq:effectivegamma} and \eqref{eq:Acapacity} imply
\begin{align}
\label{eq:gammacapacity}
\gamma_i
=HT\exp(\beta_{A0})Capacity_i^{1+\beta_{A1}}.
\end{align}
Writing the effective-coefficient specification as
$\log\gamma_i=\alpha_{\gamma0}+\alpha_{\gamma1}\log Capacity_i$ gives the
parameter mapping
\begin{align}
\label{eq:parameter_mapping}
\alpha_{\gamma0}=\log(HT)+\beta_{A0},
\qquad
\alpha_{\gamma1}=1+\beta_{A1}.
\end{align}
Thus, $\beta_{A1}=0$ corresponds to homogeneous absolute risk aversion over
total project NPV, while the elasticity of the effective coefficient
$\gamma_i$ with respect to capacity is one. A negative (positive)
$\beta_{A1}$ means that larger projects have lower (higher) absolute risk
aversion over total NPV. This parameterization separates preference
heterogeneity from the mechanical increase in total exposure associated with
project scale. In particular, a restricted model with homogeneous absolute
risk aversion over total NPV imposes $\alpha_{\gamma1}=1$, rather than
$\alpha_{\gamma1}=0$, in the effective per-unit parameterization.

I assume an exogenous wholesale price process and a common normally
distributed belief for the summary of future wholesale prices,
$r\sim\mathcal{N}(\mu_r,\sigma_r^2)$. CARA utility and normality imply that the
certainty equivalent of winning is
\begin{align}
\label{eq:ce}
CE_i^{\mathrm{total}}(q_i,p_i)
=K_i\left\{
q_i\widetilde{\delta}p_i+(1-q_i)\mu_r-c_i
-RP_r^{\mathrm{unit}}(1-q_i\mid\gamma_i)
\right\},
\end{align}
where the risk premium per unit of lifetime production is
\begin{align}
\label{eq:riskpremiumunit}
RP_r^{\mathrm{unit}}(1-q_i\mid\gamma_i)
=(1-q_i)^2\frac{\gamma_i\sigma_r^2}{2}.
\end{align}
Equivalently, the total project risk premium is
\begin{align*}
RP_r^{\mathrm{total}}(1-q_i\mid A_i,K_i)
=K_iRP_r^{\mathrm{unit}}(1-q_i\mid\gamma_i)
=(1-q_i)^2\frac{A_iK_i^2\sigma_r^2}{2}.
\end{align*}
The total risk premium increases quadratically with project scale holding
$A_i$ fixed because the standard deviation of total wholesale-market revenue
scales with $K_i$. Expressed per unit, the relevant coefficient is
$\gamma_i=A_iK_i$. Bidders earn a certainty equivalent of zero if they lose
the auction.\footnote{
\citet{Ryan2026} also assumes that bidders earn zero profit when they lose in
long-term purchase agreement auctions. I can allow a loser to earn a total
certainty equivalent $\Pi_{0i}>0$ that does not depend on the bid
$(q_i,b_i)$. The framework then identifies per-unit cost inclusive of the
reservation value, $c_i+\Pi_{0i}/K_i$, but does not separately identify
$c_i$ and $\Pi_{0i}$. Policy implications remain unchanged as long as a
counterfactual does not affect these two components differently.
}

Before the auction, bidders form a common belief about future wholesale
prices. Upon participating, bidders independently draw private types
$(A_i,c_i,Capacity_i)$ from a publicly known joint distribution, with
$A_i>0$, $c_i\in[\underline{c},\overline{c}]$, and
$Capacity_i\in[\underline{C},\overline{C}]$. Bidders observe $N$ and a
publicly known distribution of $D$ before bidding. They submit their bids,
$D$ realizes, and the auction determines purchase-agreement prices according
to the auction format. I next characterize equilibrium strategies in
sealed-bid and descending-clock auctions.



\subsection{Sealed-Bid Auctions}
\label{sec:sb}

Winners receive purchase agreements at their bid prices, $p_i=b_i$. Bidder
$i$ chooses $(q,b)$ to maximize
\begin{align}
\label{eq:sbobjective-total}
W_i(b)\,
U_i\!\left(
K_i\left\{q\widetilde{\delta}b+(1-q)\mu_r-c_i
-RP_r^{\mathrm{unit}}(1-q\mid\gamma_i)\right\}
\,\middle|\,A_i
\right),
\end{align}
where
\begin{align}
\label{eq:winprob}
W_i(b)
\equiv
\Pr\!\left(
\sum_{j\neq i}q_jCapacity_j\I(b_j\leq b)<D
\right)
\end{align}
is his winning probability. Equation~\eqref{eq:utilityequivalence} shows that
the objective in equation~\eqref{eq:sbobjective-total} is $K_i$ times the
per-unit objective indexed by $\gamma_i$. Since $K_i$ is fixed and positive,
this multiplicative term does not affect the optimal bid.

I assume that the winning probability is strictly between zero and one for
all possible bid prices. A pure-strategy Bayes Nash equilibrium (BNE),
$\{(q_i^\ast,b_i^\ast)\}_{i=1}^N$, maximizes expected utility for every bidder
under equilibrium winning-probability functions $W_i^\ast(b)$. I derive a
sufficient condition for the existence of a monotone pure-strategy BNE by
applying \citet{Reny2011}. Under the same condition, a symmetric monotone
pure-strategy BNE exists when bidders are ex-ante symmetric and independently
draw $(A_i,c_i,Capacity_i)$ from a common distribution.

Bidder $i$'s optimal share satisfies
\begin{align*}
q_i^\ast
=\min\!\left\{
\max\!\left\{\underline{q},q^{\ast\ast}(b_i^\ast\mid\gamma_i)\right\},1
\right\},
\end{align*}
where the unconstrained optimal share is
\begin{align}
\label{eq:optqsb}
q^{\ast\ast}(b\mid\gamma_i)
=1-\frac{\mu_r-\widetilde{\delta}b}{\gamma_i\sigma_r^2}.
\end{align}
The bidder chooses a 100\% share when the discounted bid price is at least as
high as the expected wholesale price. As the bid price falls below that
threshold, he shifts production toward the wholesale market because the
expected price gain compensates for additional risk. The kink therefore
identifies $\mu_r$. The interior share-price relationship has slope
$\widetilde{\delta}/(\gamma_i\sigma_r^2)$, which identifies the effective
coefficient $\gamma_i$ and, given observed $K_i$, the primitive total-NPV
coefficient $A_i=\gamma_i/K_i$.\footnote{
If the bidder's risk-premium function is continuously differentiable,
strictly increasing, and strictly convex in wholesale-market exposure over
$[0,1-\underline{q}]$, the optimal share maintains the same monotone shape.
This more general result clarifies that the share-price relationship identifies
the marginal risk premium over observed wholesale-market exposures.
}

% The existing portfolio-choice figure can be inserted here unchanged, with
% its caption revised to say that the interior slope identifies the effective
% coefficient gamma_i and hence A_i given observed K_i.

Bidders choose shares as if they care only about expected utility conditional
on winning because their shares do not affect their winning probabilities.
Consequently, $W_i^\ast$ and $c_i$ do not enter the optimal share decision
conditional on $b_i^\ast$. This separation remains crucial for identifying
risk preferences from portfolio choices before recovering costs from bid
prices.

The optimal bid price satisfies the same per-unit equation as in the original
parameterization:
\begin{align}
\label{eq:optbsb}
\underbrace{q_i^\ast\widetilde{\delta}b_i^\ast
+(1-q_i^\ast)\mu_r}_{\text{Expected revenue per unit}}
&=c_i
+\underbrace{(1-q_i^\ast)^2
\frac{\gamma_i\sigma_r^2}{2}}_{\text{Risk premium per unit}} \notag\\
&\quad+
\underbrace{
\frac{1}{\gamma_i}
\ln\!\left(
-\frac{\gamma_iq_i^\ast\widetilde{\delta}W_i^\ast(b_i^\ast)}
{dW_i^\ast(b_i^\ast)/db}+1
\right)
}_{\text{Markup per unit}}.
\end{align}
To see why cost inversion is unchanged, substitute
$\gamma_i=A_iK_i$ into the first-order condition obtained from
equation~\eqref{eq:sbobjective-total} and divide total revenues, costs, risk
premiums, and markups by $K_i$. No additional capacity term remains because
$K_i$ is fixed when bidder $i$ chooses $(q_i,b_i)$. The equation converges to
the standard risk-neutral procurement-auction inversion as $A_i\to0$, or
equivalently $\gamma_i\to0$.

\subsection{Descending-Clock Auctions}
\label{sec:dc}

I analyze descending-clock auctions in line with \citet{Milgrom1982}'s English
auction model. Institutionally, these auctions are discrete descending-clock,
pay-as-bid procedures. For tractability, I treat them as the continuous-clock
limit in which the outcome is equivalent to a uniform-price auction. Bidders
choose a cutoff---the lowest price they are willing to accept---as their bid
price $b_i$, and winners are paid a common clearing price
$p=b_{M+1:N}$, where $M$ is the number of winners.\footnote{
In the actual discrete-clock implementation, winners' final submitted prices
may not coincide because bids adjust in discrete steps with a minimum
decrement. I examine robustness to the continuous-clock approximation by
varying the equilibrium clearing-price distribution around values implied by
winners' observed final submitted prices.
}

Bidder $i$ chooses $(q,b)$ to maximize
\begin{align}
\label{eq:dcobjective-total}
\int
\I(b<p)\,
U_i\!\left(
K_i\left\{q\widetilde{\delta}p+(1-q)\mu_r-c_i
-RP_r^{\mathrm{unit}}(1-q\mid\gamma_i)\right\}
\,\middle|\,A_i
\right)f_p(p)\,dp.
\end{align}
The total-NPV objective is $K_i$ times its normalized per-unit counterpart, so
the multiplicative scale again does not affect the optimal share or cutoff.
I assume that the equilibrium clearing price
$p^\ast=b_{M+1:N}^\ast$ is normally distributed,
$\mathcal{N}(\mu_{p^\ast},\sigma_{p^\ast}^2)$, and independent of $r$. Under
the conditions detailed in the Online Appendix, a monotone pure-strategy BNE
exists, and a symmetric monotone equilibrium exists for ex-ante symmetric
bidders.

The optimal share is the constrained version of the solution to
\begin{align}
\label{eq:focqdc}
q^{\ast\ast}
=\left(
\frac{1}{1+\widetilde{\delta}^2\sigma_{p^\ast}^2/\sigma_r^2}
\right)
\left(
1-
\frac{\mu_r-
\widetilde{\delta}\widetilde{\mu}_{p^\ast}
(q^{\ast\ast},b\mid\gamma_i)}
{\gamma_i\sigma_r^2}
\right),
\end{align}
where
\begin{align*}
\widetilde{\mu}_{p^\ast}(q,b\mid\gamma_i)
=\mu_{p^\ast}
+\sigma_{p^\ast}\Lambda\!\left(
\frac{\mu_{p^\ast}-b}{\sigma_{p^\ast}}
-q\gamma_i\widetilde{\delta}\sigma_{p^\ast}
\right)
\end{align*}
is the effective expected clearing price and $\Lambda$ is the inverse Mills
ratio. The coefficient entering both expressions is
$\gamma_i=A_iK_i$. Hence, the descending-clock root problem and numerical
integration are identical to those under the per-unit formulation after this
substitution. The optimal cutoff remains the price at which expected utility
conditional on winning is zero because
$U_i(0\mid A_i)=0$ and multiplication by $K_i>0$ does not change that
condition.

The qualitative implications are therefore unchanged. The equilibrium share
increases continuously with the cutoff price, although the relationship is
nonlinear. Greater clearing-price variance makes the purchase agreement more
risky but also raises the value of truncation conditional on winning. The
effect of $A_i$ is governed by its effective counterpart $\gamma_i=A_iK_i$
because both purchase-agreement and wholesale-market risks scale with total
project production. The reparameterization changes the economic
interpretation of heterogeneity across capacities, but it does not change the
descending-clock solution for a given value of $\gamma_i$.

\bibliographystyle{aer}
\bibliography{paper/uncertainty}

\end{document}
